% =================================================================
\section{Experimental Setup}
\label{sec:experiments}
% =================================================================

\paragraph{Evaluation protocol.}
We use Leave-One-Domain-Out~(LODO): train on four of the five
recording conditions, evaluate on the fifth.
We run folds D1--D4 and report the mean; D5 is excluded because
holding it out leaves only 1,308 training clips---a data-starvation
problem~(Section~\ref{sec:data}).
All ablations use seed~42.
We separately run seeds 1234 and 3407 on the baseline and proposed
system to estimate variance; these are in Table~\ref{tab:multiseed}.

\begin{table}[h]
  \centering
  \caption{LODO D1--D4 mean $\BAu$ across seeds for baseline and
           proposed system. $^\star$D3 fold incomplete; 3-fold mean.}
  \label{tab:multiseed}
  \setlength{\tabcolsep}{5pt}
  \resizebox{\columnwidth}{!}{%
  \begin{tabular}{lccc}
    \toprule
    System & seed 42 & seed 1234 & seed 3407 \\
    \midrule
    MTRCNN baseline  & 0.166 & 0.197 & 0.147 \\
    Proposed (Bal+DANN+DiCL+proj128+$\tau$0.2) & 0.378 & $0.351^\star$ & 0.355 \\
    \bottomrule
  \end{tabular}%
  }
\end{table}

\paragraph{Metrics.}
\emph{Balanced accuracy}~(BA) is the mean per-class recall over the 9
species---it gives equal weight to rare species instead of letting the
numerically dominant species dominate the score.
Random chance for 9 classes is 11.1\%.
Following the challenge definition:
\begin{itemize}
  \item $\BAu$~(primary): BA on the held-out condition.
  \item $\BAs$: BA on the four seen conditions.
  \item $\DSG = |\BAs - \BAu|$: how much worse the model is on new
    conditions vs.\ known ones (lower is better).
\end{itemize}
All reported values are LODO D1--D4 mean.

\paragraph{Training.}
Adam optimiser, learning rate 1e-3, weight decay 1e-4, batch size 64.
We use early stopping with a minimum of 10 epochs, patience of~5
epochs, and a maximum of 100; the checkpoint with highest $\BAu$ on
the LODO validation split is used for evaluation.
Features are normalised by the per-bin mean and standard deviation
of the training split.
For DANN: $\lambda_{\max}{=}1.0$, $E{=}100$ (Eq.~\ref{eq:dann_schedule}).
For DiCL: weight $w_c{=}0.1$, temperature $\tau{=}0.2$, projection
head $32{\to}128{\to}128$ with ReLU between layers.
For FBS-Mix: $\alpha{=}0.1$, probability $p{=}0.5$, cutoff $K{=}9$.
All runs use one NVIDIA A100 GPU; each LODO fold takes roughly 27~min
for MTRCNN and 100~min for AST.

\paragraph{What we compare.}
We build up the system one component at a time to show the
contribution of each piece:
\begin{itemize}
  \item \textbf{Baseline}: unbalanced MTRCNN, standard cross-entropy.
  \item \textbf{+MixStyle}~/~\textbf{+DANN}: apply each to the
    \emph{unbalanced} data, to show they are near-useless without
    balance.
  \item \textbf{Balanced}: balanced resampling only.
  \item \textbf{Balanced + Species-only}$^*$: balanced, domain head
    turned off---tests whether domain supervision helps at all.
  \item \textbf{Balanced + DiCL}$^\dagger$: contrastive loss without
    adversarial training---tests whether DiCL works alone.
  \item \textbf{Balanced + DANN}: adversarial training on balanced data.
  \item \textbf{Balanced + DANN + DiCL}: add contrastive loss (32-dim,
    $\tau{=}0.07$, no projection head).
  \item \textbf{+proj128} / \textbf{+$\tau$0.2}: ablate the projection
    head and temperature.
  \item \textbf{Balanced + DANN + FBS-Mix}: test FBS-Mix as a
    DANN alternative (no DiCL).
  \item \textbf{Best + Augmentations$^\P$}: full proposed system plus
    HPSS and clip normalisation; results pending.
\end{itemize}
We additionally report two negative results in-text: test-time
batch-norm adaptation~(TTBN) and an Audio Spectrogram
Transformer~(AST) backbone.
